% 第6章 6.7.2：七点图上RatioCut的松弛下界、扫描候选与全部63个非空二分（由 generate.py 生成）。
% 最小值 J*=7/12≈0.5833，最大值 98/25；谱下界 lambda_2≈0.31738。
% 每个圆点是一个二分，相同取值竖直堆叠，取值相差不足 1.5mm 的两列略微错开；蓝色实心为阈值扫描的五个候选，空心为其余二分。
% 各取值及个数：0.5833×1，0.7000×1，0.9333×1，1.1667×1，1.4000×2，1.5167×1，1.7500×1，1.8200×4，2.1000×6，2.3333×9，2.5200×3，2.6833×7，2.8000×3，2.9167×6，3.0333×1，3.2200×2，3.2667×8，3.5000×5，3.9200×1
\begin{tikzpicture}[x=1mm,y=1mm,font=\figurefont\footnotesize,text=FigureInk,
  other/.style={circle,draw=FigureMuted,fill=white,line width=.5pt,minimum size=1.25mm,inner sep=0pt},
  cand/.style={circle,draw=FigureBlue,fill=FigureBlue,line width=.5pt,minimum size=1.25mm,inner sep=0pt},
  lab/.style={font=\figurefont\fontsize{7.5}{9}\selectfont},
  note/.style={font=\FigureNoteFont\fontsize{7.5}{9}\selectfont,text=FigureMuted}]
\draw[draw=FigureStroke,line width=.6pt,-{Stealth[length=1.4mm]}] (0,0)--(100.80,0);
\draw[draw=FigureStroke,line width=.6pt] (0.00,0)--(0.00,-1) node[lab,below] {$0$};
\draw[draw=FigureStroke,line width=.6pt] (24.00,0)--(24.00,-1) node[lab,below] {$1$};
\draw[draw=FigureStroke,line width=.6pt] (48.00,0)--(48.00,-1) node[lab,below] {$2$};
\draw[draw=FigureStroke,line width=.6pt] (72.00,0)--(72.00,-1) node[lab,below] {$3$};
\draw[draw=FigureStroke,line width=.6pt] (96.00,0)--(96.00,-1) node[lab,below] {$4$};
\node[lab,anchor=north east] at (100.80,-4.2) {RatioCut};
\node[cand] at (14.00,1.50) {};
\node[cand] at (16.80,1.50) {};
\node[other] at (22.40,1.50) {};
\node[cand] at (28.00,1.50) {};
\node[cand] at (33.60,1.50) {};
\node[cand] at (33.60,3.05) {};
\node[other] at (36.40,1.50) {};
\node[other] at (42.00,1.50) {};
\node[other] at (43.68,1.50) {};
\node[other] at (43.68,3.05) {};
\node[other] at (43.68,4.60) {};
\node[other] at (43.68,6.15) {};
\node[other] at (50.40,1.50) {};
\node[other] at (50.40,3.05) {};
\node[other] at (50.40,4.60) {};
\node[other] at (50.40,6.15) {};
\node[other] at (50.40,7.70) {};
\node[other] at (50.40,9.25) {};
\node[other] at (56.00,1.50) {};
\node[other] at (56.00,3.05) {};
\node[other] at (56.00,4.60) {};
\node[other] at (56.00,6.15) {};
\node[other] at (56.00,7.70) {};
\node[other] at (56.00,9.25) {};
\node[other] at (56.00,10.80) {};
\node[other] at (56.00,12.35) {};
\node[other] at (56.00,13.90) {};
\node[other] at (60.48,1.50) {};
\node[other] at (60.48,3.05) {};
\node[other] at (60.48,4.60) {};
\node[other] at (64.40,1.50) {};
\node[other] at (64.40,3.05) {};
\node[other] at (64.40,4.60) {};
\node[other] at (64.40,6.15) {};
\node[other] at (64.40,7.70) {};
\node[other] at (64.40,9.25) {};
\node[other] at (64.40,10.80) {};
\node[other] at (67.20,1.50) {};
\node[other] at (67.20,3.05) {};
\node[other] at (67.20,4.60) {};
\node[other] at (70.00,1.50) {};
\node[other] at (70.00,3.05) {};
\node[other] at (70.00,4.60) {};
\node[other] at (70.00,6.15) {};
\node[other] at (70.00,7.70) {};
\node[other] at (70.00,9.25) {};
\node[other] at (72.80,1.50) {};
\node[other] at (77.09,1.50) {};
\node[other] at (77.09,3.05) {};
\node[other] at (78.59,1.50) {};
\node[other] at (78.59,3.05) {};
\node[other] at (78.59,4.60) {};
\node[other] at (78.59,6.15) {};
\node[other] at (78.59,7.70) {};
\node[other] at (78.59,9.25) {};
\node[other] at (78.59,10.80) {};
\node[other] at (78.59,12.35) {};
\node[other] at (84.00,1.50) {};
\node[other] at (84.00,3.05) {};
\node[other] at (84.00,4.60) {};
\node[other] at (84.00,6.15) {};
\node[other] at (84.00,7.70) {};
\node[other] at (94.08,1.50) {};
\draw[draw=FigureRed,dashed,line width=.8pt] (7.62,-1)--(7.62,24);
\node[lab,text=FigureRed,anchor=south west] at (5.62,24) {谱下界$\lambda_2\approx0.317$};
\draw[draw=FigureBlue,line width=.6pt,-{Stealth[length=1.2mm]}] (20.00,12)--(14.60,2.6);
\node[lab,anchor=south west,text=FigureBlue] at (18.50,12) {$J^\star=U=\tfrac7{12}$};
\node[cand] at (60,24) {}; \node[note,anchor=west] at (61.5,24) {阈值扫描的五个候选};
\node[other] at (60,19.5) {}; \node[note,anchor=west] at (61.5,19.5) {其余二分};
\end{tikzpicture}
